\chapter{Conclusion}\doublespacing
\label{Chapter6}
\lhead{Chapter VI.}

This thesis set out to answer the question: \emph{among a
representative set of credit-risk classifiers, which model performs
most consistently across multiple available UCI datasets, and is
its advantage statistically significant relative to the others?}
The answer that emerges from the empirical work of
Chapters~\ref{Chapter4} and~\ref{Chapter5} is:

\begin{quote}
\textbf{The Stacked Ensemble, with its meta-learner chosen by
$F_1$-macro on out-of-fold predictions of five diverse base learners,
is the most consistently strong classifier across the four UCI
credit datasets. Its advantage is supported by a Friedman test
that rejects the global null at $p<0.001$ on every metric and by a
Wilcoxon-Holm pairwise post-hoc that establishes
\emph{Stacker $>$ XGBoost} at $\alpha=0.05$ on AUC.}
\end{quote}

\section{Summary of Findings}
\label{sec:summary}

The headline numbers are gathered in Table~\ref{tab:summary} for ready reference.

\begin{table}[H]
\centering
\caption{Summary of the cross-dataset and statistical results.
Stacker is the best classifier on every aggregate metric and the
test statistic rejects the global null at $p<0.001$ throughout.}
\singlespacing
\label{tab:summary}
\small
\renewcommand{\arraystretch}{1.2}
\begin{tabular}{l c c c}
\toprule
\textbf{Result} & \textbf{$F_1$-macro} & \textbf{AUC} & \textbf{G-Mean}\\
\midrule
Stacker mean rank (low\,=\,best) & \textbf{2.00} & \textbf{1.38} & \textbf{2.19}\\
Best non-stacked model & XGBoost (2.75) & RF (2.59) & XGBoost (3.66)\\
Friedman $\chi^{2}_F$ (dof = 7) & 46.20 & 76.69 & 34.07\\
Friedman $p$-value & $<0.001$ & $<0.001$ & $<0.001$\\
Iman--Davenport $F_F$ & 10.53 & 32.58 & 6.56\\
Holm-significant pairs ($\alpha=0.05$) & 8 of 28 & 16 of 28 & 7 of 28\\
\hspace{4mm}of which \emph{Stacker $>$ X} & 3 & \textbf{6} & 5\\
\bottomrule
\end{tabular}
\end{table}
The empirical contributions can be summarised in four points:

\begin{enumerate}
 \item \textbf{Descriptive ranking.} Across the 16 evaluation
 blocks the Stacked Ensemble has the lowest average rank on
 $F_1$-macro (2.00), AUC (1.38) and G-Mean (2.19). Gradient
 Boosting is the lowest-rank single (non-stacked) classifier
 on accuracy. The classical baselines (Logistic Regression,
 KNN, Decision Tree, SVM-RBF) consistently fill the lower half
 of the ranking.
 \item \textbf{Global significance.} The Friedman test rejects
 the global null hypothesis of equal classifier performance at
 $p<0.001$ on every one of the four metrics ($F_1$-macro, AUC,
 accuracy, G-Mean). The Iman--Davenport $F$-correction agrees.
 Whatever metric the credit-risk analyst chooses, the eight
 classifiers do not perform equally on these data.
 \item \textbf{Pairwise significance.} The Wilcoxon-Holm
 post-hoc resolves \emph{Stacker $>$ XGBoost} at $\alpha=0.05$
 on AUC ($p=0.0008$) and at $\alpha=0.10$ on G-Mean
 ($p=0.0535$). The Stacker is significantly better than six
 of seven competitors on AUC and four of seven on G-Mean.
 On $F_1$-macro the leading trio (Stacker, XGBoost, Gradient
 Boosting) is statistically tight; the Stacker is the
 descriptive but not the inferential winner on this metric.
 \item \textbf{Methodological by-products.} The GA-MaxAccLogit
 meta-learner --- an adaptation of the ProfLogit
 \citep{stripling2018proflogit} genetic algorithm in which the
 expected-profit fitness is replaced by $F_1$-macro --- is
 selected as the winning meta on a non-trivial fraction of the
 (dataset, FS) pairs (most notably on the German and Taiwan
 datasets). The diverse winning metas (Logistic Regression,
 GA-MaxAccLogit, MLP, XGBoost) across the 16 blocks justify
 the meta-candidate sweep: a fixed meta would not have
 extracted the same performance.
\end{enumerate}

\section{Contributions Revisited}
\label{sec:contributions-revisited}

In the language of Chapter~\ref{Chapter1}, the four contributions
have all been delivered:
\begin{itemize}
 \item A \emph{uniform pipeline} that applies the same
 preprocessing, feature-selection sweep, hyperparameter-tuning
 protocol and evaluation suite to all eight classifiers on all
 four UCI datasets (Chapters~\ref{Chapter3} and~\ref{Chapter4}).
 \item The \emph{GA-MaxAccLogit} meta-learner adapted from
 ProfLogit (Section~\ref{sec:gamal}).
 \item A \emph{statistical-significance layer} on the 16-block
 rank matrix (Sections~\ref{sec:friedman}--\ref{sec:posthoc}).
 \item A \emph{reproducible Colab notebook bundle} containing
 the Friedman test, the Wilcoxon-Holm pairwise post-hoc, the
 GA-MaxAccLogit implementation and the data-loading scripts
 (Appendix~\ref{AppendixA}).
\end{itemize}

\section{Why the Stacked Ensemble is the Recommended Model}
\label{sec:why-stacker}

Before listing the limitations, this section states explicitly ---
and defends --- the headline recommendation: \emph{the Stacked
Ensemble is the model a lender should adopt}.  The argument follows
the structure used in the credit-scoring literature that motivates
the thesis: pair a descriptive ranking with a non-parametric
statistical test, and accept the model that wins on both.

\textbf{Descriptive evidence.}  Across the sixteen evaluation
blocks (4 datasets $\times$ 4 feature-selection variants), the
Stacker has the lowest average rank on every aggregate metric:
2.00 on $F_1$-macro, 1.38 on AUC, and 2.19 on G-Mean.  On AUC the
Stacker is the per-dataset best on \emph{every one} of the four
datasets, a sweep that no competitor achieves.  It is never lower
than rank \#2 on any individual block --- something we cannot say
of XGBoost (rank \#3 on Taiwan) or Gradient Boosting (rank \#4 on
Taiwan).

\textbf{Inferential evidence.}  The Friedman rank test rejects the
global null ``the eight classifiers perform equally'' at $p<0.001$
on every one of the four metrics (Table~\ref{tab:friedman}).  The
Iman--Davenport $F$-correction agrees.  The post-hoc Wilcoxon
signed-rank test with Holm correction then resolves the pairwise
verdicts: \emph{Stacker $>$ XGBoost} is significant at
$\alpha=0.05$ on AUC ($p=0.0008$), and the Stacker dominates 6 of
7 pairs on AUC and 5 of 7 on G-Mean.  Whatever metric the lender
weights most, the Stacker is either the statistically significant
winner (AUC, G-Mean) or the descriptive winner ($F_1$-macro,
accuracy).

\textbf{Why not the runners-up.}  XGBoost is the strongest single
non-stacked model: it has the highest single-row $F_1$-macro on
two of four datasets (Australian Lasso-LR 0.876; German RFE 0.708).
But XGBoost is volatile across feature-selection variants ---
its German mean $F_1$-macro across the four FS variants drops to
0.655 because three of the four FS variants are weak.  A lender
cannot rerun a feature-selection sweep on every batch of
applicants; they pin one pipeline and run it.  The Stacker's
narrower spread across FS variants makes it the lower-risk
deployment choice.

Gradient Boosting is the steadiest \emph{single} model (mean rank
2.75 on $F_1$-macro), but it is significantly worse than the
Stacker on AUC ($p<0.001$ Wilcoxon-Holm) and trails by 0.014 on
mean AUC.  Random Forest, although strong on AUC (mean 0.851), is
weaker on the threshold-dependent metrics because of its
conservative leaf voting on imbalanced data.

\textbf{Why $F_1$-macro is the headline classification metric.}
The lender's two errors --- approving a defaulter (FN) and
rejecting a creditworthy applicant (FP) --- have very different
business consequences, but the dataset itself does not know which
loss matters more without a cost matrix that the bank must supply.
In the absence of a fixed cost matrix, the metric that is fair to
\emph{both} classes is $F_1$-macro: the unweighted mean of the
per-class $F_1$ scores.  Plain $F_1$ on the positive class can be
inflated by a model that sacrifices the minority for the majority
--- on Taiwan (78\%/22\%) a trivial majority-vote classifier
scores 0.78 accuracy but $F_1=0$ on the default class.
$F_1$-macro punishes that asymmetry, while AUC measures the
threshold-independent ranking power.  Together they give the
lender a complete picture of how the classifier behaves both at
the chosen operating point ($F_1$-macro) and across all operating
points (AUC).

\textbf{Why AUC carries the burden of the formal claim.}  GHOST
threshold tuning is fitted on the training set and is
non-deterministic across runs; consequently
\emph{threshold-dependent} metrics (accuracy, $F_1$-macro, F-Score,
Cohen's $\kappa$, G-Mean, Youden's J) wobble by approximately
$\pm 0.02$ between runs on German and Taiwan.  AUC, by contrast,
integrates over the full ROC curve and therefore depends only on
the classifier's ranking, not on the chosen threshold; it is
\emph{identical} run-to-run on every classifier in our
experiments.  This is why every formal claim
above rests on AUC: it is the one metric for which the post-hoc
verdict is not contaminated by GHOST stochasticity.

\section{Limitations}
\label{sec:limitations}

We close the thesis by stating the limitations of the present
work, in roughly increasing order of importance.

\textbf{Number of datasets.} Four datasets is at the low end of
what \citet{demvsar2006statistical} recommends for the Friedman
rank test.  We compensate by treating each (dataset, FS) pair as
an independent block, which raises the effective $n$ to 16, but
the four-dataset constraint is real: results may shift if
additional credit datasets (Polish, Loan Prediction, Kaggle
``Give Me Some Credit'') are added.

\textbf{GHOST threshold stochasticity.}  The GHOST threshold-tuning
step involves a random search whose seed is not fully controlled.
Consequently the threshold-dependent metrics on German and Taiwan
wobble by approximately $\pm 0.02$ across runs.  AUC, by contrast,
is threshold-independent and bit-for-bit reproducible across runs;
this is why the formal Wilcoxon-Holm claims in
Section~\ref{sec:why-stacker} rest on AUC rather than $F_1$-macro
on those two datasets.  The descriptive rank ordering of
classifiers is unaffected.

\textbf{Single seed for the train/test split.} We use a fixed
random seed (42) for the 75/25 stratified train/test split.  A
multi-seed evaluation would give confidence intervals on the
per-classifier metrics; we do not report these.

\section{Future Work}
\label{sec:future}

The thesis is built around two classification metrics ---
$F_1$-macro and AUC --- and the diagnostic ten-measure suite that
surrounds them.  These are the right objectives for a study whose
goal is to compare classifiers on a level playing field, but they
are not the objectives a deployed credit-risk model is ultimately
judged by.  The natural direction for future work is therefore to
\emph{move from classification accuracy to expected profit}.

In a deployed lending decision, the two errors carry very different
monetary consequences: a false negative (approving a defaulter) loses
the loan principal, while a false positive (rejecting a creditworthy
applicant) loses only the future interest revenue.  Maximising
$F_1$-macro treats these two errors symmetrically; maximising
expected profit does not.  Restoring the original ProfLogit fitness
function --- expected profit minus the L1 penalty --- inside the
GA-MaxAccLogit meta-learner is therefore the natural next step.
Concretely, Equation~\eqref{eq:gamal}'s fitness becomes
\[
\mathrm{fitness}(\boldsymbol{\theta})=\mathrm{EMP}(\boldsymbol{\theta})\;-\;\lambda\lVert\boldsymbol{\beta}\rVert_1,
\]
where $\mathrm{EMP}(\boldsymbol{\theta})$ is the expected maximum
profit for credit scoring computed at the classifier's decision
threshold.  Because the rest of the GA machinery
(selection, crossover, mutation, elitism, soft-thresholding) is
gradient-free, the only change required is in the fitness function.
The Stacker would then optimise not the metric of \emph{statistical
correctness} but the metric the lender's profit-and-loss statement
cares about.

A second, more straightforward extension is to validate the
present findings on a broader collection of credit datasets.  The
patterns observed here --- the Stacker's robustness, the GHOST
threshold sensitivity, the dominance of repayment-status features
for default prediction --- should hold on any sufficiently similar
benchmark; testing on additional datasets would confirm the
generality of the conclusion and would also increase the power of
the Friedman and Wilcoxon-Holm post-hoc tests so that comparisons
the current four-dataset benchmark could not resolve at
$\alpha=0.05$ become decidable.

With these two extensions the present benchmark would become a
profit-aware, broadly validated credit-risk study --- the natural
next step beyond this thesis.
